\subsection{Sprint 1 Results}
\label{sec:results-sprint1}

The CI pipeline did not exist during Sprint 1.
It was created on 2026-11-04, in Sprint 2.
For this reason the same suites were run on 2026-11-20 on the Sprint 1 code, with the commands of the current pipeline from Section~\ref{sec:ci}.

The code tested is the deployed one: tag \texttt{sprint1-deploy}, commit \sha{3bbef9d}.
It is the merge of pull request \#44 into \texttt{main} (\sha{ca1799a}, 2026-10-04) plus one deployment-configuration commit.
It predates the first organization pull request, so it holds only Sprint 1 work.
Evidence files are \evid{EV-S1-01} to \evid{EV-S1-08}.

\begin{table}[H]
  \centering
  \begin{tabularx}{\linewidth}{L{3cm}L{2.6cm}L{1.4cm}L{1.4cm}L{1.3cm}L{1.4cm}Y}
    \toprule
    Suite & Level & Tests & Passed & Failed & Skipped & Evidence \\
    \midrule
    Backend, Unit & Unit & 151 & 151 & 0 & 0 & \evid{EV-S1-01} \\
    Backend, Feature & Integration & 700 & 700 & 0 & 0 & \evid{EV-S1-01} \\
    Frontend & Unit, component & 518 & 518 & 0 & 0 & \evid{EV-S1-02} \\
    Leodeguita & Unit, component & 157 & 157 & 0 & 0 & \evid{EV-S1-03} \\
    \midrule
    \textbf{Total} & & \textbf{1526} & \textbf{1526} & \textbf{0} & \textbf{0} & \\
    \bottomrule
  \end{tabularx}
  \caption{Sprint 1 automated tests, final results after the re-test described below.}
  \label{tab:results-sprint1-automated}
\end{table}

The backend uses PHPUnit 11.5 and the two frontends use Vitest.
Coverage was not measured, because that step was added to the pipeline later.

\evidencefigure{tests/sprint1-backend-phpunit.png}{Backend, PHPUnit result. Tag \texttt{sprint1-deploy}, \sha{3bbef9d}, run on 2026-11-20.}{fig:results-sprint1-backend}
\evidencefigure{tests/sprint1-frontend-vitest.png}{Frontend, Vitest result. Tag \texttt{sprint1-deploy}, \sha{3bbef9d}, run on 2026-11-20.}{fig:results-sprint1-frontend}
\evidencefigure{tests/sprint1-leodeguita-vitest.png}{Leodeguita, Vitest result. Tag \texttt{sprint1-deploy}, \sha{3bbef9d}, run on 2026-11-20.}{fig:results-sprint1-leodeguita}

\subsubsection*{High-priority test cases}

Case IDs come from the test plan in Section~\ref{sec:test-plan}.
Integration cases are PHPUnit tests, and their results come from the report \evid{EV-S1-01}.
Functional cases run on the deployed system.

\begin{longtable}{@{}L{1.9cm}L{1.9cm}L{3.4cm}L{1.7cm}L{1.4cm}L{1.1cm}L{1.9cm}@{}}
  \toprule
  Case ID & Level & Expected & Actual & Result & Defect & Evidence \\
  \midrule
  \endhead
  \multicolumn{7}{@{}l}{\textit{HUA-01, fire department permit validation (High)}} \\
  TC-HUA-01-01 & Integration & Approved, moderation record written, manager notified & As expected & Pass & --- & \evid{EV-S1-01} \\
  TC-HUA-01-02 & Integration & Rejection with a reason code is recorded & As expected & Pass & --- & \evid{EV-S1-01} \\
  TC-HUA-01-03 & Integration & Publication request without permit is blocked & As expected & Pass & --- & \evid{EV-S1-01} \\
  TC-HUA-01-04 & Integration & Rejection without a reason code is refused & As expected & Pass & --- & \evid{EV-S1-01} \\
  TC-HUA-01-05 & Integration & Reason ``otro'' without a comment is refused & As expected & Pass & --- & \evid{EV-S1-01} \\
  TC-HUA-01-06 & Integration & No permit, approval without waiver is refused & As expected & Pass & --- & \evid{EV-S1-01} \\
  TC-HUA-01-07 & Integration & No permit, approval with waiver records the timestamp & As expected & Pass & --- & \evid{EV-S1-01} \\
  \midrule
  \multicolumn{7}{@{}l}{\textit{HUG-04, warehouse registration with permit (High)}} \\
  TC-HUG-04-01 & Integration & Registration with permit creates the warehouse as pending verification & As expected & Pass & --- & \evid{EV-S1-01} \\
  TC-HUG-04-02 & Integration & Submission without permit is blocked & As expected & Pass & --- & \evid{EV-S1-01} \\
  TC-HUG-04-03 & Integration & Submission with incomplete fields is blocked & As expected & Pass & --- & \evid{EV-S1-01} \\
  \bottomrule
  \caption{Sprint 1 test cases of High-priority stories.}
  \label{tab:results-sprint1}
\end{longtable}

All 10 High-priority cases passed.
They are verified by PHPUnit integration tests, which call the API and check the database.
No test drives the web interface, so these cases are not end-to-end tests.
Table~\ref{tab:cases-sprint1} lists the test behind each case that is not in the figure below.
Cases of Medium and Low priority follow the prioritization rule of Section~\ref{sec:test-plan}.

\begin{table}[H]
  \centering
  \begin{tabularx}{\linewidth}{L{2.6cm}Y}
    \toprule
    Case ID & PHPUnit test \\
    \midrule
    TC-HUG-04-01 & \texttt{StoreRoomTest::test\_\allowbreak valid\_\allowbreak permit\_\allowbreak end\_\allowbreak to\_\allowbreak end\_\allowbreak and\_\allowbreak bracket\_\allowbreak notation\_\allowbreak prices} \\
    TC-HUG-04-02 & \texttt{StoreRoomTest::test\_\allowbreak permit\_\allowbreak missing\_\allowbreak returns\_\allowbreak 400\_\allowbreak with\_\allowbreak exact\_\allowbreak message} \\
    TC-HUG-04-03 & \texttt{StoreRoomTest::test\_\allowbreak create\_\allowbreak store\_\allowbreak room\_\allowbreak without\_\allowbreak store\_\allowbreak prices\_\allowbreak is\_\allowbreak rejected}, \texttt{StoreRoomTest::test\_\allowbreak missing\_\allowbreak cancellation\_\allowbreak policy\_\allowbreak tier\_\allowbreak returns\_\allowbreak 400} \\
    TC-HUA-01-03 & Same test as TC-HUG-04-02, seen from the administrator story. \\
    \bottomrule
  \end{tabularx}
  \caption{Tests behind the Sprint 1 cases not shown in the figure.}
  \label{tab:cases-sprint1}
\end{table}

\evidencefigure{tests/sprint1-backend-hua01.png}{HUA-01, the six integration cases of the table above (\texttt{TC-HUA-01-01, 02, 04, 05, 06, 07}). Tag \texttt{sprint1-deploy}, \sha{3bbef9d}, run on 2026-11-20.}{fig:results-sprint1-hua01}

\subsubsection*{Summary}

All 1526 tests passed in the final run, with none skipped.
No defects were recorded.

\subsubsection*{Failures and Re-tests}

Two failures appeared, and both came from the environment, not from the application.
\begin{itemize}
  \item On the first attempt, all 157 Leodeguita tests failed because Node 26 has a native \texttt{localStorage} that the tests do not expect.
  After disabling that feature, they passed on the same commit.
  \item On the run at \sha{3bbef9d}, one Frontend test, \texttt{LeodegaUI overlap message and server errors (RB-4)}, failed with a timeout of 5000\,ms.
  The machine was slow in that run, and the suite took 72\,s instead of 36\,s.
  The complete suite was run again and all 518 tests passed.
  The report of the failed run is kept as \evid{EV-S1-04}.
\end{itemize}
